\chapter{Code Generation and Runtime Support}
\label{ch:codegen-runtime}

After IR generation, the compiler emits C code for the selected field and target. This chapter describes the generated API, byte-to-field mapping, data structures, load/store routines, arithmetic, carry propagation, and reduction loops.

\section{Generated C API}
\label{sec:generated-c-api}

Every DSL hash function is emitted as a C function with the signature
\begin{lstlisting}[style=thesisC]
void h(uint8_t *output, uint8_t *key,
       uint8_t *message, size_t len);
\end{lstlisting}
The parameter \(len\) is a byte length at the C boundary. The generated function converts it to a field-element count using \(c=\lfloor \pi/8\rfloor\).

\paragraph{Buffer convention.}
The caller must provide input buffers with enough padding for generated loads. The compiler computes the block count by ceiling division and may read slightly past the logical end.

\paragraph{Output format.}
The output buffer is assumed large enough for the exported field element. Stores write the final reduced value in little-endian byte order. The number of exported bytes is \(\lceil \pi/8\rceil\), while input chunks consume \(\lfloor \pi/8\rfloor\) bytes.

\section{Generated Files and Build Integration}
\label{sec:generated-files}

For a module \(m\), the compiler emits \(m\texttt{.h}\), \(m\texttt{.c}\), and shared data-structure files. It can also emit a C++ test file and a C benchmarking source.

The surrounding CMake project builds correctness and performance targets and links handwritten tests with dependencies such as OpenSSL~\cite{openssl} or Boost~\cite{boost}.

\section{Big Integer Representation}
\label{sec:bigint-representation}

Generated C represents a field element as a struct with one entry per limb. Scalar field elements use machine-word limbs; vectorized field elements use SIMD vector types.

The data structures also define field-derived constants: limb count, limb mask, most-significant-limb size, input chunk size, and exported byte count.

\Cref{lst:generated-structs} shows a generated data-structure fragment for \(2^{116}-3\) on two-lane NEON, with \(\lambda=18\), \(\lambda'=8\), and \(\ell=7\).

\begin{lstlisting}[style=thesisC,basicstyle=\scriptsize\ttfamily,caption={Generated scalar and vectorized field-element data structures.},label={lst:generated-structs}]
#include <arm_neon.h>

#define FIELD_PI 116ull
#define FIELD_THETA 3ull
#define FIELD_LAMBDA 18ull
#define FIELD_LAMBDA_PRIME 8ull
#define FIELD_LIMBS 7ull
#define FIELD_CHUNK_SIZE 14ull
#define FIELD_EXPORT_BYTES 15ull
#define FIELD_PRIME_HEX "0xffffffffffffffffffffffffffffd"

typedef struct {
    uint64_t limb0;
    uint64_t limb1;
    uint64_t limb2;
    uint64_t limb3;
    uint64_t limb4;
    uint64_t limb5;
    uint64_t limb6;
} bigint_t;

typedef struct {
    uint64x2_t limb0;
    uint64x2_t limb1;
    uint64x2_t limb2;
    uint64x2_t limb3;
    uint64x2_t limb4;
    uint64x2_t limb5;
    uint64x2_t limb6;
} vbigint_t;
\end{lstlisting}

\section{Load and Store}
\label{sec:load-store}

Loading is field-specific because it maps byte chunks into limbs. Scalar load reads consecutive bytes and splits them into limbs. Vectorized load assembles several field elements at once, one per SIMD lane.

Loaded input field elements contain only \(\lfloor \pi/8\rfloor\) bytes. If \(\pi\) is not divisible by eight, the remaining high bits are not filled by consuming another byte.

Stores perform the reverse operation after full reduction. The store is part of the IR, so hash functions do not need a wrapper to export bytes.

\section{Generated Arithmetic}
\label{sec:generated-arithmetic}

Addition is emitted as component-wise limb addition.

\paragraph{Multiplication and squaring.}
Multiplication is emitted as schoolbook or Karatsuba. Schoolbook computes limb products and folds high terms using the Crandall identity. Karatsuba reduces limb multiplications at the cost of additions and register pressure.

Squaring has a dedicated implementation that computes cross terms once and doubles them.

For multiplication, the generated code follows the folded schoolbook formula from \Cref{sec:prime-field-arithmetic}. Each output limb accumulates ordinary products plus wrapped products multiplied by the Crandall factor.

\Cref{lst:squaresum-neon} continues the square-sum example from \Cref{lst:squaresum-dsl,lst:squaresum-ir}. The excerpt is generated NEON C for \(GF(2^{116}-3)\), showing the field-element count, vector loop, part of \texttt{vsquare}, horizontal reduction, scalar tail, and final reduction.

\begin{lstlisting}[style=thesisC,basicstyle=\scriptsize\ttfamily,caption={Excerpt of generated NEON C for the square-sum running example.},label={lst:squaresum-neon}]
void squaresum(uint8_t* output, uint8_t* key, uint8_t* message, size_t buffer_length) {
const uint64x2_t vlambda_mask = vdupq_n_u64(262143);
const uint64x2_t vlambda_prime_mask = vdupq_n_u64(255);
const uint64_t len = (buffer_length + FIELD_CHUNK_SIZE - 1) / FIELD_CHUNK_SIZE;
uint64_t var0 = len-0;
uint64_t var1 = var0%2;
uint64_t var2 = len-var1;
vbigint_t var3 = { 0 };
for (uint64_t i=0;i<var2;i+=2) {
/* var4 = vload message @ i offsets 0, 1 */
vbigint_t var4;
uint8x16x2_t var4_limb0_data = {
vld1q_u8(&message[14 * (i + 0) + 0]),
vld1q_u8(&message[14 * (i + 1) + 0]),
};
var4.limb0 =
vandq_u64(vreinterpretq_u64_u8(vqtbl2q_u8(var4_limb0_data, (uint8x16_t){
    0, 1, 2, 255, 255, 255, 255, 255,
    16, 17, 18, 255, 255, 255, 255, 255,
})),vlambda_mask);
/* remaining vload limbs omitted */

vbigint_t var5 = {0};
/* var5 = vsquare var4 */
var5.limb0 = vmlal_u32(
    var5.limb0,
    vmovn_u64(var4.limb0),
    vmovn_u64(var4.limb0)
);
var5.limb0 = vmlal_u32(
    var5.limb0,
    vmul_u32(vmovn_u64(var4.limb1), vdup_n_u32(3072)),
    vmovn_u64(var4.limb6)
);
/* remaining vsquare terms omitted */

var3.limb0 = vaddq_u64(var3.limb0, var5.limb0);
/* remaining accumulator limbs omitted */
}
/* var6 = horiz_add var3 */
bigint_t var6;
var6.limb0 = vaddvq_u64(var3.limb0);
/* remaining horizontal-add limbs omitted */

for (uint64_t i=var2;i<len;i+=1) {
/* var7 = load message @ i offsets 0 */
bigint_t var7;
var7.limb0 = (message[14 * i + 0] << 0|message[14 * i + 1] << 8|message[14 * i + 2] << 16) & 262143ull;
/* remaining scalar load limbs omitted */
bigint_t var8 = {0};
/* var8 = square var7 */
var8.limb0 += var7.limb0 * var7.limb0;
var8.limb0 += 6144 * var7.limb1 * var7.limb6;
/* remaining scalar square and add terms omitted */
}
/* var6 = full_reduction var6 */
/* full_reduction: extended carry */
var6.limb1 += var6.limb0 >> 18;
var6.limb0 &= 262143ull;var6.limb2 += var6.limb1 >> 18;
/* remaining reduction and store bytes omitted */
}
\end{lstlisting}

\section{Carry and Full Reduction}
\label{sec:carry-full-reduction}

Generated carry code propagates overflow bits from one limb to the next. The carry out of the most significant limb is multiplied by \(\theta\) and folded into the least significant limb. Vectorized code performs the same operation lane-wise.

A full reduction is emitted at the end of hash functions. It handles cases where carry propagation alone leaves a value above the prime.

\section{Backends}
\label{sec:backends}

The scalar backend is useful in its own right and provides a baseline for vectorization speedups.

\paragraph{Vector backends.}
The NEON backend targets two 64-bit lanes. The AVX2 backend targets four lanes with the same limb-major layout, but different loads, lane-power generation, and loop structure.

\section{Performance Harness}
\label{sec:performance-harness}

The compiler can emit a benchmark program that allocates key, message, and output buffers, calls the generated hash over a range of byte lengths, and reports cycles per call. The harness treats the hash function as a black box depending only on the C API.
